\section{Feature-Cut Priorities}
\label{sec:feature-cuts}

This program is not optimized around producing the maximum number of boards. It is optimized around reaching November~30 with measured evidence that LAN8660~+~LAN8661 make the Greenfield Alfred E2B architecture viable, and MCU~+~LAN8651 makes the Brownfield Alfred Gateway architecture viable, with 10BASE-T1S as the common wedge.

\subsection{Protect at all costs}

\begin{enumerate}[itemsep=0.1em]
\item Known-good T1S reference network (Section~\ref{sec:cots-bench}).
\item LAN8660 validation (Section~\ref{sec:lan8660-validation}).
\item LAN8661 validation (Section~\ref{sec:lan8661-validation}).
\item PLCA validation (Section~\ref{sec:t1s-plca}).
\item Rev A.
\item Rev B.
\item Greenfield proof (Section~\ref{sec:arch-a-proof}).
\item CAN Brownfield Gateway (Section~\ref{sec:arch-b}).
\item Passive mode.
\item Authorized bidirectional mode.
\item Quantitative T1S validation (Section~\ref{sec:measurement}).
\end{enumerate}

\subsection{Cut first}

\begin{itemize}[itemsep=0.1em]
\item LIN.
\item 100BASE-T1 / 1000BASE-T1.
\item Large node counts (beyond the mandatory minimum plus the preferred second LAN8660).
\item Additional peripheral types beyond the mandatory motor-driver/actuator and lighting paths.
\item Sophisticated enclosure.
\item Additional Gateway protocols.
\item Cosmetic features.
\item Rev~C, if necessary.
\end{itemize}

\begin{decision}[title={LAN8660 and LAN8661 are explicitly not on the cut list}]
They are protected at the same priority as the T1S reference network itself, because cutting either one removes the entire Greenfield proof, not a feature of it. If schedule pressure threatens them specifically, the response is to cut everything in the ``cut first'' list before touching either part's validation scope, and, if truly necessary, to reduce the Greenfield demo to the single-mandatory-LAN8660 configuration (Section~\ref{sec:arch-a}) rather than dropping LAN8660 or LAN8661 altogether.
\end{decision}
